\documentclass[11pt]{article}
\usepackage[margin=1in]{geometry}
\usepackage{amsmath}
\usepackage{hyperref}
\usepackage{enumitem}

\title{Clearing Simulation: Single-Run Technical Report and Onboarding Guide}
\author{Project Clearing Simulation}
\date{\today}

\begin{document}
\maketitle

\section{Purpose and Scope}
This report explains how a single simulation run works in the clearing project,
focusing on the logic, algorithms, and daily processing steps. It references the
codebase locations that implement each component. The second half provides
user guidelines for onboarding new team members, including a step-by-step daily
timeline of the simulation phases.

\section{High-Level Architecture}
\subsection{Core Modules}
\begin{itemize}[leftmargin=*]
  \item \texttt{scripts/run\_simulation.py}: CLI entry point. Loads RBM model and quantizer,
  initializes portfolios, sets simulation parameters, calls \texttt{simulate\_day}, prints
  per-day progress and end-of-run summary.
  \item \texttt{clearing/simulation.py}: Main daily loop in \texttt{simulate\_day}. Handles
  settlement, margining, defaulting, trading proposal, QUBO construction, and acceptance.
  \item \texttt{clearing/market.py}: Market state transitions and return generation.
  \item \texttt{clearing/sampler.py}: Scenario sampling via RBM with clamped visible units.
  \item \texttt{clearing/quantization.py}: Return quantization and bit encoding for the RBM.
  \item \texttt{clearing/margin.py}: Expected Shortfall (ES) and VaR computation.
  \item \texttt{clearing/trading.py}: Trade proposal model.
  \item \texttt{clearing/qubo.py}: Builds the BQM for trade acceptance and solves it.
  \item \texttt{clearing/logger.py}: Logging of daily phases and optional scenario stats.
\end{itemize}

\subsection{Key Data Shapes}
Throughout the simulation:
\begin{itemize}[leftmargin=*]
  \item $P \in \mathbb{R}^{M \times N}$: client portfolios (M clients, N instruments).
  \item $W \in \mathbb{R}^M$: total wealth per client.
  \item $C \in \mathbb{R}^M$: collateral per client.
  \item $r_t \in \mathbb{R}^N$: daily market returns.
  \item $R \in \mathbb{R}^{\Omega \times N}$: scenario return matrix.
  \item \texttt{alive} $\in \{\texttt{True}, \texttt{False}\}^M$: whether client remains active.
\end{itemize}
These shape checks are enforced in \texttt{clearing/simulation.py::simulate\_day}.

\section{Market Model}
The market model uses a Markov chain with three regimes (calm, volatile, crisis)
and factor-driven returns. Implemented in \texttt{clearing/market.py}.

\subsection{State Transition}
State $z_t$ is sampled using a row of transition matrix $P_z$ (see
\texttt{make\_default\_params}). The default transition matrix biases toward staying in the
current regime.

\subsection{Return Generation}
Given $z_t$, factor returns $f_t$ and idiosyncratic noise $\epsilon_t$ are sampled:
\[
f_t \sim \mathcal{N}(0, \sigma_f[z_t]), \quad
\epsilon_t \sim \mathcal{N}(0, \sigma_e[z_t])
\]
Returns are:
\[
r_t = B f_t + \epsilon_t
\]
If \texttt{rho} is nonzero, returns mix with $r_{t-1}$ using
$r_t = \rho r_{t-1} + (1-\rho) r_{\text{raw}}$.
Optional jumps are applied with probability \texttt{jump\_p[z\_t]} using
\texttt{jump\_sigma} and \texttt{jump\_mu}. See \texttt{clearing/market.py::generate\_day}.

\section{Scenario Generation via RBM}
Scenario sampling is conditional on the current regime $z_t$ and previous returns
$r_{t-1}$. Implemented in \texttt{clearing/sampler.py::sample\_scenarios\_from\_float}.

\subsection{Quantization}
Returns are quantized per instrument into $K$ bits using quantile-based bounds
\texttt{ReturnQuantizer} in \texttt{clearing/quantization.py}.
The pipeline:
\begin{enumerate}[leftmargin=*]
  \item Compute per-instrument quantile bounds \texttt{fit\_return\_quantizer}.
  \item Encode $r_{t-1}$ into $K$-bit vectors (\texttt{encode\_returns}).
  \item Clamp visible units for $z_t$ (one-hot) and $r_{t-1}$ in the RBM.
  \item Sample the RBM with Gibbs sampling (\texttt{RBM.sample\_clamped}).
  \item Decode sampled bits into scenario returns (\texttt{decode\_returns}).
\end{enumerate}
The RBM configuration (visible block sizes) is loaded from the model run folder
in \texttt{scripts/run\_simulation.py}.

\section{Margin Model (ES/VaR)}
Margins are computed from scenario losses using Expected Shortfall (ES):
\begin{align*}
\text{PnL} &= R P^T \\
\text{Loss} &= -\text{PnL}
\end{align*}
Losses are sorted per client, and ES is the mean of the worst $(1-\alpha)$ tail.
VaR is the loss at quantile $\alpha$. Implemented in
\texttt{clearing/margin.py::margin\_es}.

\section{Daily Simulation Algorithm}
The daily loop is in \texttt{clearing/simulation.py::simulate\_day}. The phases below
match the logger phases used by the Streamlit dashboard.

\subsection{Phase 1: Start of Day}
Inputs: $P$, $W$, $C$, \texttt{alive}, previous regime $z_{t-1}$ and returns $r_{t-1}$.
Logged as phase \texttt{start}.

\subsection{Phase 2: Market Move}
Generate market regime and returns with \texttt{generate\_day}. Compute PnL:
\[
\text{PnL}_i = \sum_{j=1}^N P_{ij} \cdot r_{t,j}
\]
Logged as phase \texttt{market\_move}.
If stress is enabled, returns are adjusted on a chosen stress day:
\[
r_t \leftarrow r_t \cdot \texttt{stress\_return\_scale} + \texttt{stress\_return\_shift}
\]
See \texttt{SimDayParams} in \texttt{clearing/simulation.py}.

\subsection{Phase 3: Settlement}
Wealth and collateral are updated:
\begin{align*}
W_{\text{settle}} &= W + \text{PnL} \\
F_{\text{pre}} &= (W - C) + \text{PnL} \\
C_{\text{after}} &= C + \min(F_{\text{pre}}, 0)
\end{align*}
Defaults are accounted for as:
\[
\text{default\_loss\_step} = \max(-C_{\text{after}}, 0)
\]
Collateral and wealth are floored at zero. Total default loss accumulates in
\texttt{default\_loss}. Logged as phase \texttt{settlement}.

\subsection{Phase 4: Margining}
Scenario returns $R$ are sampled from the RBM and used to compute ES and VaR.
Required margin $M_{\text{req}}$ is the ES per client. Logged as phase \texttt{margin}.

\subsection{Phase 5: Default and Liquidity}
Clients must top up collateral to $M_{\text{req}}$ using available funds:
\begin{align*}
F_{\text{avail}} &= \max(W_{\text{settle}} - C_{\text{after}}, 0) \\
\text{need} &= \max(M_{\text{req}} - C_{\text{after}}, 0)
\end{align*}
If $F_{\text{avail}} < \text{need}$, the client defaults. Defaulted clients are
removed from trading, and their portfolios are liquidated to zero
(\texttt{P\_liq}). Logged as phase \texttt{default}.

\subsection{Phase 6: Trade Proposal}
Proposed trades are generated for active clients in
\texttt{clearing/trading.py::propose\_trades}:
\begin{align*}
\mu &= \texttt{base\_scale} \cdot \texttt{momentum\_weight} \cdot r_t \cdot (|P|+1)^{\gamma} \\
\sigma &= \texttt{noise\_scale} \cdot (\sqrt{|P|} + 1) \\
\Delta P &= \mu + \sigma \cdot \epsilon,\quad \epsilon \sim \mathcal{N}(0, I)
\end{align*}
Zero positions are traded with probability \texttt{p\_trade\_when\_zero}, optionally
with a separate noise term \texttt{zero\_trade\_std}. Trades are clipped by
\texttt{max\_abs\_trade}. Logged as phase \texttt{trades}.

\subsection{Phase 7: QUBO Clearing}
Tentative margins are computed for $P_{\text{tent}} = P_{\text{liq}} + \Delta P$.
Margin change is:
\[
\Delta M = M_{\text{req,tent}} - M_{\text{req,cur}}
\]
A binary decision vector $x$ selects which clients' trades to accept.
The BQM is built in \texttt{clearing/qubo.py::build\_bqm\_accept\_clients}:
\begin{itemize}[leftmargin=*]
  \item Penalizes positive $\Delta M$ to keep total risk below budget $B$.
  \item Rewards negative $\Delta M$ (risk-reducing trades).
  \item Optional utility term is based on trade size
  (\texttt{utility\_weight} times $|\Delta P|$).
\end{itemize}
The solver is either simulated annealing (\texttt{sa}) or hybrid (\texttt{hybrid}),
implemented in \texttt{clearing/qubo.py::solve\_bqm}. Hybrid uses
\texttt{DWAVE\_API\_KEY} and supports labels such as \texttt{ClSim\{run\}\_day\{d\}}
via the \texttt{label} parameter. Logged as phase \texttt{decision}.

\subsection{Phase 8: End of Day}
Accepted trades update portfolios:
\[
P_{\text{next}} = P_{\text{liq}} + x \odot \Delta P
\]
Margins are recomputed for logging. Outputs include updated $P$, $W$, $C$,
\texttt{alive}, default loss totals, QUBO stats, and updated state $z_t$, $r_t$.
Logged as phase \texttt{end}.

\section{Initialization and Run Loop}
\texttt{scripts/run\_simulation.py} performs the following:
\begin{itemize}[leftmargin=*]
  \item Loads an RBM run folder and optional quantizer. If the quantizer is missing,
  it is fitted from a fresh market simulation.
  \item Initializes portfolios $P$ with random values and a density mask.
  \item Computes initial margin and sets collateral and liquidity buffers.
  \item Runs \texttt{simulate\_day} for $T$ days, prints day-by-day progress.
  \item Prints a final summary: Total DL, Clients' gain, total defaults, mean QUBO
  size/time/energy, and total simulation time.
\end{itemize}
Stress-test presets in \texttt{run\_simulation.py} can amplify volatility and jump
sizes, increase initial position sizes, reduce buffers, and trigger a stress day.

\section{User Guidelines for New Members}
\subsection{Typical Workflow}
\begin{enumerate}[leftmargin=*]
  \item Install dependencies from \texttt{requirements.txt}.
  \item (Optional) Generate a dataset and quantizer:
  \texttt{python scripts/generate\_data.py --T 10000 --out-csv data/data.csv --out-quantizer data/quantizer.pt}
  \item Run a single simulation:
  \texttt{python scripts/run\_simulation.py --model-run models/model\_mini --quantizer data/quantizer.pt}
  \item View results in the Streamlit dashboard:
  \texttt{streamlit run streamlit\_app.py} and open the generated \texttt{run\_log.pt}.
\end{enumerate}

\subsection{Daily Timeline (What Happens Each Day)}
Use this checklist when reading logs or onboarding:
\begin{enumerate}[leftmargin=*]
  \item \textbf{Start:} Capture $P$, $W$, $C$, and the previous market state.
  \item \textbf{Market move:} Sample $z_t$ and $r_t$, compute PnL.
  \item \textbf{Settlement:} Update wealth/collateral and accumulate default loss.
  \item \textbf{Margins:} Sample scenarios and compute ES/VaR.
  \item \textbf{Default:} Check funding shortfalls and deactivate defaulted clients.
  \item \textbf{Trades:} Propose trades based on momentum and exposure.
  \item \textbf{Decision:} Solve the QUBO to accept or reject trades.
  \item \textbf{End:} Apply accepted trades and log final state for the day.
\end{enumerate}

\subsection{Interpreting End-of-Run Summary}
At the end of a simulation, the console prints:
\begin{itemize}[leftmargin=*]
  \item \textbf{Total DL:} Sum of default losses over all days.
  \item \textbf{Clients' gain:} Total wealth change from start to end.
  \item \textbf{Total defaulted clients:} Count of clients that could not meet margin.
  \item \textbf{Mean QUBO size/time/energy:} Average solver statistics per day.
  \item \textbf{Total simulation time:} Wall-clock runtime for the full run.
\end{itemize}

\subsection{Benchmarking}
The benchmarking tool \texttt{benchmark.py} runs multiple simulations with fixed
settings, then exports histograms and metadata. The output files are named
\texttt{benchmark\_\{solver\}\_\{date\_time\}.png} and
\texttt{benchmark\_\{solver\}\_\{date\_time\}.json}.

\section{Practical Notes}
\begin{itemize}[leftmargin=*]
  \item Scenario generation assumes RBM visible blocks match $N \cdot K$.
  If you change $N$ or quantization bits $K$, you must retrain or regenerate
  the RBM and quantizer.
  \item The parameter \texttt{post\_full\_margin\_daily} is currently implemented
  but uses the same logic in both branches; it exists for future extensions.
  \item Hybrid QUBO requires \texttt{DWAVE\_API\_KEY} to be set in the environment.
\end{itemize}

\end{document}
